\documentclass[a4paper,12pt]{article}

\usepackage{geometry}
\geometry{top=20mm,bottom=20mm,left=25mm,right=25mm}

\usepackage[utf8]{inputenc}
\usepackage[T2A]{fontenc}
\usepackage[russian]{babel}
\usepackage{amsmath}

\usepackage{xcolor}
\usepackage{listings}
\usepackage{hyperref}

\usepackage{booktabs}
\usepackage{array}

\definecolor{codebg}{RGB}{248,248,248}
\definecolor{codeframe}{RGB}{210,210,210}
\definecolor{codekeyword}{RGB}{0,76,153}
\definecolor{codecomment}{RGB}{0,128,0}
\definecolor{codestring}{RGB}{163,21,21}
\definecolor{codenumber}{RGB}{120,120,120}

\lstdefinestyle{common}{
    backgroundcolor=\color{codebg},
    frame=single,
    rulecolor=\color{codeframe},
    basicstyle=\ttfamily\small,
    keywordstyle=\color{codekeyword}\bfseries,
    commentstyle=\color{codecomment}\itshape,
    stringstyle=\color{codestring},
    numberstyle=\tiny\color{codenumber},
    numbers=left,
    numbersep=10pt,
    breaklines=true,
    showstringspaces=false,
    keepspaces=true,
    columns=fullflexible,
    tabsize=4
}

\lstdefinestyle{cppstyle}{
    style=common,
    language=C++
}

\lstdefinestyle{pythonstyle}{
    style=common,
    language=Python
}

\begin{document}

\begin{titlepage}
    \centering
    {\large МИНОБРНАУКИ РОССИИ\par}
    \vspace{0.5em}
    {\large САНКТ-ПЕТЕРБУРГСКИЙ ГОСУДАРСТВЕННЫЙ\par}
    {\large ЭЛЕКТРОТЕХНИЧЕСКИЙ УНИВЕРСИТЕТ\par}
    {\large «ЛЭТИ» ИМ. В.И. УЛЬЯНОВА (ЛЕНИНА)\par}
    \vspace{1em}
    {\large Кафедра Информационных систем\par}

    \vspace{6em}

    {\LARGE \textbf{Практическая работа}\par}
    \vspace{0.8em}
    {\large по дисциплине «Теория принятия решений»\par}
    \vspace{0.8em}
    {\large Тема: «Применение методов линейного программирования для решения задачи поиска оптимального плана заказа стали на завод»\par}

    \vfill

    \begin{flushleft}
        Студент гр. 3376 \hfill Михайлов Н.Ф.\\[0.5em]
        Преподаватель   \hfill Пономарёв А.В.
    \end{flushleft}

    \vspace{3em}

    Санкт-Петербург\par
    2026
\end{titlepage}

\tableofcontents
\newpage

\section{Постановка задачи}

Рассматривается задача перевозки стали двух марок $a$ и $b$
со складов $A_1$, $A_2$ в пункты потребления $B_1$, $B_2$.

Известны:
\begin{itemize}
    \item запасы стали на складах;
    \item потребности пунктов;
    \item стоимость перевозки;
    \item возможность замещения стали марки $a$ в $b$.
\end{itemize}

Требуется найти план перевозок, минимизирующий суммарные затраты, и провести исследование оптимального решения.

\section{Формализация задачи}

\subsection{Формальная постановка}

\textbf{Переменные.}

\begin{itemize}
    \item $x_{ij}^a$, где $i,j \in \{1,2\}$ --- количество стали марки $a$,
    перевозимой со склада $A_i$ в пункт $B_j$, т;
    \item $x_{ij}^b$, где $i,j \in \{1,2\}$ --- количество стали марки $b$,
    перевозимой со склада $A_i$ в пункт $B_j$, т;
    \item $y_{ij}$, где $i,j \in \{1,2\}$ --- объем стали марки $a$,
    направляемый по маршруту $A_i \to B_j$ на замену стали марки $b$, т.
\end{itemize}

Переменная $y_{ij}$ не описывает отдельную перевозку и поэтому не входит в
целевую функцию. Она только отражает использование уже доставленной стали
марки $a$ вместо стали марки $b$. При этом $1$ т стали марки $a$ заменяет
$0.7$ т стали марки $b$.

Тогда математическая модель имеет вид:

\begin{equation}
\min Z =
2(x_{11}^a + x_{21}^a + x_{11}^b + x_{21}^b) +
5.3(x_{12}^a + x_{22}^a) +
c_b(x_{12}^b + x_{22}^b)
\end{equation}

\begin{equation}
x_{11}^a + x_{12}^a \le 5100
\end{equation}

\begin{equation}
x_{21}^a + x_{22}^a \le 2900
\end{equation}

\begin{equation}
x_{11}^b + x_{12}^b \le 3200
\end{equation}

\begin{equation}
x_{21}^b + x_{22}^b \le 5600
\end{equation}

\begin{equation}
x_{11}^a + x_{21}^a - y_{11} - y_{21} \ge 2400
\end{equation}

\begin{equation}
x_{11}^b + x_{21}^b + 0.7y_{11} + 0.7y_{21} \ge 4500
\end{equation}

\begin{equation}
x_{11}^a + x_{21}^a + x_{11}^b + x_{21}^b \ge 8300
\end{equation}

\begin{equation}
x_{12}^a + x_{22}^a - y_{12} - y_{22} \ge 3400
\end{equation}

\begin{equation}
x_{12}^b + x_{22}^b + 0.7y_{12} + 0.7y_{22} \ge 2300
\end{equation}

\begin{equation}
x_{12}^a + x_{22}^a + x_{12}^b + x_{22}^b \ge 6300
\end{equation}

\begin{equation}
y_{11} \le x_{11}^a,\quad
y_{12} \le x_{12}^a,\quad
y_{21} \le x_{21}^a,\quad
y_{22} \le x_{22}^a
\end{equation}

\begin{equation}
x_{ij}^a \ge 0,\quad x_{ij}^b \ge 0,\quad y_{ij} \ge 0
\end{equation}

Для базового решения и анализа правых частей принималось $c_b = 5.3$.
При анализе целевой функции параметр $c_b$ изменялся как стоимость
перевозки стали марки $b$ по направлению в пункт $B_2$.

\subsection{Табличное представление}

Для перехода к форме $Gx \le h$ порядок переменных выбирался следующим:
\[
\left(
x_{11}^a,\,
x_{12}^a,\,
x_{21}^a,\,
x_{22}^a,\,
x_{11}^b,\,
x_{12}^b,\,
x_{21}^b,\,
x_{22}^b,\,
y_{11},\,
y_{12},\,
y_{21},\,
y_{22}
\right)
\]

\begin{center}
\small
\setlength{\tabcolsep}{4pt}
\begin{tabular}{>{\centering\arraybackslash}m{1.7cm}cccccccccccccc}
\toprule
Огр. &
$x_{11}^a$ & $x_{12}^a$ & $x_{21}^a$ & $x_{22}^a$ &
$x_{11}^b$ & $x_{12}^b$ & $x_{21}^b$ & $x_{22}^b$ &
$y_{11}$ & $y_{12}$ & $y_{21}$ & $y_{22}$ &
Знак & Пр. часть \\
\midrule
$c$ & 2 & 5.3 & 2 & 5.3 & 2 & $c_b$ & 2 & $c_b$ & 0 & 0 & 0 & 0 & min & -- \\
$g_1$  & 1 & 1 & 0 & 0 & 0 & 0 & 0 & 0 & 0 & 0 & 0 & 0 & $\le$ & 5100 \\
$g_2$  & 0 & 0 & 1 & 1 & 0 & 0 & 0 & 0 & 0 & 0 & 0 & 0 & $\le$ & 2900 \\
$g_3$  & 0 & 0 & 0 & 0 & 1 & 1 & 0 & 0 & 0 & 0 & 0 & 0 & $\le$ & 3200 \\
$g_4$  & 0 & 0 & 0 & 0 & 0 & 0 & 1 & 1 & 0 & 0 & 0 & 0 & $\le$ & 5600 \\
$g_5$  & -1 & 0 & -1 & 0 & 0 & 0 & 0 & 0 & 1 & 0 & 1 & 0 & $\le$ & -2400 \\
$g_6$  & 0 & 0 & 0 & 0 & -1 & 0 & -1 & 0 & -0.7 & 0 & -0.7 & 0 & $\le$ & -4500 \\
$g_7$  & -1 & 0 & -1 & 0 & -1 & 0 & -1 & 0 & 0 & 0 & 0 & 0 & $\le$ & -8300 \\
$g_8$  & 0 & -1 & 0 & -1 & 0 & 0 & 0 & 0 & 0 & 1 & 0 & 1 & $\le$ & -3400 \\
$g_9$  & 0 & 0 & 0 & 0 & 0 & -1 & 0 & -1 & 0 & -0.7 & 0 & -0.7 & $\le$ & -2300 \\
$g_{10}$ & 0 & -1 & 0 & -1 & 0 & -1 & 0 & -1 & 0 & 0 & 0 & 0 & $\le$ & -6300 \\
$g_{11}$ & -1 & 0 & 0 & 0 & 0 & 0 & 0 & 0 & 1 & 0 & 0 & 0 & $\le$ & 0 \\
$g_{12}$ & 0 & -1 & 0 & 0 & 0 & 0 & 0 & 0 & 0 & 1 & 0 & 0 & $\le$ & 0 \\
$g_{13}$ & 0 & 0 & -1 & 0 & 0 & 0 & 0 & 0 & 0 & 0 & 1 & 0 & $\le$ & 0 \\
$g_{14}$ & 0 & 0 & 0 & -1 & 0 & 0 & 0 & 0 & 0 & 0 & 0 & 1 & $\le$ & 0 \\
\bottomrule
\end{tabular}
\end{center}

Неотрицательность переменных в коде задается отдельно ограничениями вида
$-x_k \le 0$ для каждой из 12 переменных.

\pagebreak
\subsection{Особенность численного решения}

Задача решалась в два этапа. Сначала минимизировались суммарные транспортные
затраты. Затем среди всех решений с той же минимальной стоимостью выбиралось
решение с минимальным суммарным замещением $\sum y_{ij}$. Такой подход нужен,
поскольку при одинаковой стоимости перевозок для обеих марок задача может иметь
несколько равноправных оптимальных планов, часть из которых содержит замещение,
а часть --- нет.

Формально процедура имела следующий вид. На первом этапе решалась задача
\[
\min c^T x \quad \text{при} \quad Gx \le h,
\]
где вектор $c$ содержит коэффициенты транспортных затрат. В результате
получалось оптимальное значение целевой функции $Z^*$.

После этого решалась вторая задача:
\[
\min \sum y_{ij}
\]
при тех же ограничениях модели и дополнительном условии
\[
c^T x \le Z^* + \varepsilon,
\]
где $\varepsilon$ --- малый численный допуск. Тем самым из множества всех
планов, оптимальных по затратам, выбирался план с минимально возможным
объемом замещения. Именно эта процедура использовалась в коде при вызове
решателя CVXOPT.

\section{Решение задачи}

Задача решалась численно с использованием библиотеки CVXOPT.

\subsection{Базовое решение}

Получен следующий оптимальный план:

\textbf{Сталь марки $a$:}
\[
\begin{aligned}
A_1 \to B_1 &: 2024.15 \\
A_1 \to B_2 &: 2216.86 \\
A_2 \to B_1 &: 1194.73 \\
A_2 \to B_2 &: 1328.55
\end{aligned}
\]

\textbf{Сталь марки $b$:}
\[
\begin{aligned}
A_1 \to B_1 &: 1823.98 \\
A_1 \to B_2 &: 1013.52 \\
A_2 \to B_1 &: 3257.15 \\
A_2 \to B_2 &: 1741.07
\end{aligned}
\]

\textbf{Замещение:}
\[
y_{11} = 0,\quad y_{12} = 0,\quad y_{21} = 0,\quad y_{22} = 0
\]

\[
\sum y_{ij} = 0
\]

\[
Z = 49990
\]

\textbf{Пояснение.}

В базовом решении замещение отсутствует. Это объясняется тем, что суммарный
запас стали марки $b$ на складах равен
\[
3200 + 5600 = 8800 \text{ т},
\]
а этого достаточно, чтобы покрыть обязательный спрос на сталь марки $b$
($4500 + 2300 = 6800$ т) и дополнительную потребность в 2000 т стали любой
марки. Поэтому оптимальный план не использует замену стали марки $a$ на $b$.

\subsection{Анализ изменения правых частей}

Так как в базовом решении замещение отсутствует, в соответствии с условием
выполнялось постепенное увеличение запасов стали марки $a$ и одновременное
уменьшение запасов стали марки $b$ на обоих складах до появления замещения
не менее 100 т.

Получены результаты:

\begin{center}
\begin{tabular}{cccc}
\toprule
Шаг & Запасы $a$ & Запасы $b$ & Замещение \\
\midrule
0 & 5100 / 2900 & 3200 / 5600 & 0 \\
10 & 5600 / 3400 & 2700 / 5100 & 0 \\
20 & 6100 / 3900 & 2200 / 4600 & 0 \\
21 & 6150 / 3950 & 2150 / 4550 & 142.86 \\
\bottomrule
\end{tabular}
\end{center}

\textbf{Вывод:}

Замещение появляется только после достижения порогового значения запасов.
До шага 20 имеющегося объема стали марки $b$ достаточно для покрытия спроса,
поэтому замещение не используется. На шаге 21 запас стали марки $b$
уменьшается настолько, что оптимальный план начинает использовать замену
стали марки $a$ на $b$, причем объем замещения составляет $142.86$ т.

\subsection{Анализ целевой функции}

Для исключения влияния ограничений по запасам сначала были существенно
увеличены запасы обеих марок: до 10000 т на каждом складе. После этого
выполнялось постепенное увеличение стоимости перевозки стали марки $b$
по направлению в пункт $B_2$.

Результаты:

\begin{itemize}
    \item при $c_b = 5.3$: замещение отсутствует
    \item при $c_b = 5.5$: замещение составляет $2000$ т
\end{itemize}

\textbf{Вывод:}

Изменение коэффициента целевой функции приводит к резкой перестройке структуры
оптимального плана. При стоимости $c_b = 5.3$ перевозка стали марки $b$
остается выгодной, и замещение не требуется. При увеличении стоимости до
$5.5$ ДЕ за тонну возникает замещение объемом 2000 т, то есть модель
переходит к использованию стали марки $a$ для покрытия части спроса на
сталь марки $b$ по направлению в пункт $B_2$.

\section{Листинг кода}

\subsection{Реализация на Python}

\lstinputlisting[style=pythonstyle]{lab2_listing.py}


\end{document}
